%NO MODIFICAR ESTA SECCION!
\documentclass[12pt]{article} % Define la clase del documento, en este caso, un artículo
\usepackage{mathptmx}
\usepackage[letterpaper,margin=3cm]{geometry} % Configura el tamaño del papel y los márgenes del documento
\usepackage{graphicx} % Permite la inserción de imágenes
\usepackage[spanish]{babel}% Activar esta configuración para informes en español, ajusta el idioma del documento
\usepackage[usenames]{color} % Permite el uso de colores definidos por nombre en el documento
\usepackage{hyperref} % Habilita enlaces y referencias dentro del documento
\usepackage{booktabs} % Proporciona comandos para crear tablas de alta calidad
\usepackage{natbib}
\usepackage{tikz} % Permite la creación de gráficos y diagramas vectoriales directamente en LaTeX
\usepackage{float} % Para controlar la posición de los elementos flotantes, como imágenes, con la opción [H]
\usepackage{diagbox} % Permite crear celdas con líneas diagonales en tablas
\usepackage{listings} % Permite la inclusión y formateo de código fuente en el documento
\usepackage{xcolor} % Paquete para definir y usar colores en el documento
\usepackage{parskip} % Añade espacio entre párrafos en lugar de sangrías
\usepackage{fancyhdr} % Permite personalizar encabezados y pies de página
\usepackage{amsmath} % Proporciona una amplia variedad de entornos y comandos matemáticos
\usepackage{microtype}  % Mejora la justificación del texto

\hypersetup{
    colorlinks   = true,    % Colorea los enlaces
    linkcolor    = blue,    % Color de los enlaces internos (por ejemplo, en la tabla de contenidos)
    citecolor    = blue,    % Color de los enlaces a citas bibliográficas
    filecolor    = blue,    % Color de los enlaces a archivos locales
    urlcolor     = blue,    % Color de los enlaces externos (URLs)
}

\geometry{
    paperwidth=21.6cm,  % Ancho del papel
    paperheight=27.9cm,  % Largo del papel
    left=3cm,  % Margen izquierdo
    right=2cm,  % Margen derecho
}

\pagestyle{fancy} % Usa el estilo fancyhdr
\fancyhf{} % Borra todos los encabezados y pies de página
\renewcommand{\headrulewidth}{0pt}
\renewcommand{\footrulewidth}{0pt} % Desactiva la línea horizontal predeterminada en el pie
%\setlength{\headheight}{1.5cm} % Ajusta la altura del encabezado para hacer espacio para la línea
%\fancyhead[L]{\raisebox{0.20cm}{\textbf{Sistemas de Transporte}}} % Añade el texto en la parte izquierda del encabezado, subiéndolo ligeramente
%\fancyhead[R]{\raisebox{0.1cm}{\includegraphics[width=0.25\linewidth]{LOGO_UNIVERSIDAD.jpg}}} % Añade la imagen en la parte derecha del encabezado y súbela un poco
%\fancyhead[C]{\rule{\textwidth}{0.6pt}} % Añade una línea horizontal superior centrada
%\fancyfoot[C]{\rule{\textwidth}{0.6pt}} % Añade una línea horizontal en el pie de página centrada
%\fancyfoot[R]{\raisebox{-1.5\baselineskip}{\thepage}} % Coloca el número de página a la derecha, con suficiente espacio debajo de la línea
%\geometry{top=2.5cm, bottom=2.5cm} % Ajusta los márgenes superior e inferior

\fancyhead[R]{\thepage} % Número de página a la derecha en el encabezado


% Definición de colores al estilo Visual Studio Code
\definecolor{codegreen}{rgb}{0.25,0.49,0.48} % Comentarios
\definecolor{codegray}{rgb}{0.5,0.5,0.5} % Números y anotaciones
\definecolor{codepurple}{rgb}{0.58,0,0.82} % Palabras clave
\definecolor{backcolour}{rgb}{0.95,0.95,0.92} % Color de fondo

% Configuración del estilo de las celdas de código
\lstset{
    backgroundcolor=\color{backcolour},   % color de fondo; necesita que el paquete color o xcolor esté cargado
    commentstyle=\color{codegreen},       % estilo de comentarios
    keywordstyle=\color{codepurple},      % estilo de palabras clave
    numberstyle=\tiny\color{codegray},    % estilo de los números de línea
    stringstyle=\color{red},              % estilo de las cadenas de texto
    basicstyle=\ttfamily\small,           % estilo del texto básico
    breakatwhitespace=false,              % ajustes de líneas sólo en espacios en blanco
    breaklines=true,                      % ajustar las líneas si son muy largas
    captionpos=b,                         % posición de la leyenda (abajo)
    keepspaces=true,                      % preserva los espacios en el texto; útil si se usa monoespaciado
    numbers=left,                         % dónde poner los números de línea
    numbersep=5pt,                        % qué tan lejos están los números de línea del código
    showspaces=false,                     % mostrar espacios con subrayados particulares; reemplaza 'showstringspaces'
    showstringspaces=false,               % subrayar los espacios dentro de las cadenas solo
    showtabs=false,                       % mostrar tabulaciones en el código con subrayados particulares
    tabsize=2,                            % tamaños de tabulación a 2 espacios
    language=TeX,                         % lenguaje del código
    morecomment=[l]\#,                    % reconocer # como inicio de comentario en Python
    frame=single,                         % agregar un marco simple alrededor del código
    rulecolor=\color{black}               % color del marco
}

\begin{document}
%----------------------------------------------------------------------------------------
%   PORTADA
%Modificar desde aqui en adelante
%----------------------------------------------------------------------------------------
\begin{titlepage}%Inicio de la carátula, solo modificar los datos necesarios
\newcommand{\HRule}{\rule{\linewidth}{0.5mm}} 
\center 
%----------------------------------------------------------------------------------------
%	ENCABEZADO
%----------------------------------------------------------------------------------------
%----------------------------------------------------------------------------------------
%	SECCION DEL TITULO
%----------------------------------------------------------------------------------------
\begin{center}
    \textbf{\LARGE Generación y Atracción de Viajes} \\[0.5cm]
    \textbf{Felipe Vicencio y Lukas Wolff} \\
    Facultad de Ingenieria y Ciencias Aplicadas, Universidad de los Andes, Santiago de Chile.\\
    email: \href{mailto:lwolff@miuandes.cl}{lwolff@miuandes.cl}, \href{mailto:favicencio@miuandes.cl}{favicencio@miuandes.cl}
    \\
    github: \href{https://github.com/LukasWolff2002/TAREA_2_AUTITOS.git}{Link al repositorio}
\end{center}

\vspace{1cm}

\begin{center}
    \textbf{\large RESUMEN}    
\end{center}

\begin{flushleft}
    Hablar aqui del resumen del informe actual, no puede exceer las 250 palabras. Poner palabras clave, ademas de selccionar 3 palabras clave por los autores
\end{flushleft}


\vspace{1cm}

\end{titlepage}
%----------------------------------------------------------------------------------------
%	SECCION DE LA FECHA
%----------------------------------------------------------------------------------------
%{\large \textbf{\today}}\\[2cm] % El comando \today coloca la fecha del dia, y esto se actualiza con cada compilacion, en caso de querer tener una fecha estatica, reemplazar el \today por la fecha deseada

%----------------------------------------------------------------------------------------
%  INDICE
%----------------------------------------------------------------------------------------
%\newpage
%\tableofcontents
%\thispagestyle{plain} % Deshabilita el encabezado en la página del índice
%\thispagestyle{empty} % Deshabilita el número de página en la página del índice
%\newpage

%Se puede agregar un indice de figuras si es nesesario
%\newpage
%\listoffigures 
%\thispagestyle{plain} % Deshabilita el encabezado en la página del índice %
%\thispagestyle{empty}
%\newpage
%----------------------------------------------------------------------------------------
%   ACÁ EMPIEZA EL INFORME
\setcounter{page}{1} % Reinicia el contador de páginas
%----------------------------------------------------------------------------------------
%Este es el formato a seguir para los titulos de las secciones

\section{Introducción}

El estudio de los viajes es fundamental para el desarrollo de las zonas altamente pobladas, ya que permite ahorrar tiempos de viaje, ahorros en combustible e incluso, mejorar el bienestar de las personas. Para poder generar estos analisis, es nesesario recaudar data real sobre el comportamiento de las personas como su tipo de transporte, origne y destino, horas de viajes, entre otros.
\\ \\
El problema de estos estudios radica en la cantidad de datos que se nesesitan para poder realizar un analisis completo, por lo tanto, se buscan hacer estimaciones mediante una regrecion lineal para poder representar a una poblacion de la manera mas eficaz posible.
\\ \\
En el presente informe, se buscara analizar el numero de viajes de la comuna de Las Condes, en base a la data obtenida de la encuesta ESI, para los años 2012, 2017 y 2023. Para esto, se utilizara un modelo de regresion lineal para poder estimar el numero de viajes de la comuna en base a distintos factores que pueden influir en la cantidad de viajes.
\\ \\
Finalmente, se analizara y discutira sobre la eficiencia del modelo, la precision de los datos obtenidos y los factores que pueden influir en la cantidad de viajes de la comuna de Las Condes.

\section{Contenido}

\section{Ecuaciones, Tablas y Figuras}

\subsection{Matriz Origen Destino}

Para la representacion de todos los viajes de la red, se puede utilizar la matriz origen destino, la cual se puede representar de la siguiente manera:

\begin{equation}
    \begin{bmatrix}
        O_{1}D_{1} & O_{1}D_{2} & \cdots & O_{1}D_{n} \\
        O_{2}D_{1} & O_{2}D_{2} & \cdots & O_{2}D_{n} \\
        \vdots & \vdots & \ddots & \vdots \\
        O_{n}D_{1} & O_{n}D_{2} & \cdots & O_{n}D_{n}
    \end{bmatrix}
\end{equation}

El problema es que normalmente no se tienen suficientes dato para poder completar de manera satisfactoria la matriz origen destino, por lo que se deben utilizar modelos de estimacion de viajes para poder completarla.
\\ \\
De esta manera, es importante analizar los distintos factores que pueden influir en los viajes de las personas.

\subsection{Análisis Regresión Lineal}

En base a los datos obtenidos por distintas encuetas, se busca obtener una relacion que explique el comportamiento de la poblacion:

\begin{equation}
    Y = \alpha + \beta_{1}X_{1} + \beta_{2}X_{2} + \cdots + \beta_{n}X_{n} + \epsilon
\end{equation}

Lo cual se puede extrapolar a:

\begin{equation}
    Y = \alpha + \vec{\beta}\vec{X} + \epsilon
\end{equation}

Donde Y corresponde al numero de viajes y $\vec{X}$ corresponden a los distintos factores que pueden afectar a la cantidad de viajes. En el caso de este informe, se utilizaran las siguientes variables X:

\begin{center}
    $X_i = log_{10}(IPCH)$, donde IPCH correspone al ingreso per capita por hogar\\
    $X_p$ = numero de personas con edad $\epsilon$ [0, 5]\\
    $X_e$ = numero de personas con edad $\epsilon$ [6, 22]\\
    $X_t$ = numero de personas con edad $\epsilon$ [23, 62]\\
    $X_j$ = numero de personas con edad $\epsilon$ [63, 79]\\
\end{center}

De esta forma, el model ode regrecion lineal queda de la siguiente forma:

\begin{equation}
    Y_i = -2.1723 + 0.3792X_i + 0.6221X_p + 1.0065X_e + 0.4302X_t + 0.1614X_j
    \label{eq:regrecion}
\end{equation}

Es importante mencionar que los coeficientes $\vec{\beta}$ y $\alpha$ fueron entregados por el enunciado, donde el modelo fue calibrado segun los datos del 2012.

\section{Resultados Bases de Datos}
DESPUES NO PONER LAS SECCIONES DE ESTA FORMA.

\subsection{Resultados 2.1}

Los datos muestrales fueron obtenidos de la encuesta ESI \textbf{\cite{esi}}

\begin{table}[H]
    \centering
    \caption{Matriz de estimacion de viajes}
    \vspace{0.2cm}
    \begin{tabular}{|c|c|c|c|c|c|c|}
        \hline
        Año & Numero Hogares & $X_i$ & $X_p$ & $X_e$ & $X_t$ & $X_j$ \\
        \hline
        2012 & 271 &5.730 & 0.258 & 0.771 & 1.978 & 0.303 \\
        2017 & 272 &5.773 & 0.243 & 0.691 & 1.978 & 0.320 \\
        2023 & 187 &5.934 & 0.160 & 0.668 & 1.775 & 0.299 \\
        \hline
    \end{tabular}
    \vspace{0.2cm}
    \\Fuente: Elaboracion propia a partir de los datos de encueta ESI
\end{table}

\subsection{Resultados 2.2}
A continuación se presentarán los resultaos obtenidos de la base de datos ``Estimación y Proyección 2002-2035``. Se agruparon las personas por grupo etario en la comuna de Las Condes para los años 2012, 2017 y 2023 y con el total de hogares se obtuvieron los siguientes parámetros.

\begin{table}[H]
    \centering
    \caption{Valor esperado de personas por grupo etario}
    \vspace{0.2cm}
    \begin{tabular}{|c|c|c|c|c|c|}
        \hline
        \diagbox{Año}{Coeficiente} & $X_p$ & $X_e$ & $X_t$ & $X_j$ & Total hogares \\
        \hline
        2012 & 0.189 & 0.608 & 1.511 & 0.417 & 104649 \\
        2017 & 0.172 & 0.515 & 1.465 & 0.453 & 118007 \\
        2023 & 0.173 & 0.509 & 1.641 & 0.567 & 118007 \\
        \hline
    \end{tabular}
    \vspace{0.2cm}
    \\Fuente: Elaboracion propia a partir de los datos de Estimación y Proyección 2002 - 2035
\end{table}

Los datos anteriores representan la cantidad de personas por rango etario presente en un hogar para cada año. 
Se realizó el análisis con datos poblacionales, es decir, con el total de hogares inscritos en el año. Se supuso la misma
cantidad de hogares para los años 2017 y 2023 debido a la falta de información proporcionada.

\subsection{Resultados 2.3}

De la regrecion lineal obtenida, se obtuvieron los siguientes resultados:

\begin{table}[H]
    \centering
    \caption{Estimacion de viajes}
    \vspace{0.2cm}
    \begin{tabular}{|c|c|c|}
        \hline
        Año & $Y_i$ & Viajes Totales \\
        \hline
        2012 & 1.84 & 192554.16\\
        2017 & 1.77 & 208872.39\\
        2023 & 1.66 & 195891.62\\
        \hline
    \end{tabular}
    \vspace{0.2cm}
    \\ Fuente: Elaboracion propia a partir de los resultados de la ecuacion \ref{eq:regrecion}
\end{table}

\section{Discusiones}

\begin{itemize}
    \item ¿Qué tan precisos cree que son cada uno de los resultados? ¿Por qué?
    \item ¿Es suficiente la información del modelo de regresión lineal obtenido?
    \item ¿Qué cosas han cambiado entre el 2012 y el 2023 que puedan haber tenido un impacto significativo en la generación de viajes, más allá de las incluidas en el modelo?
\end{itemize}

\subsection{Inflacion}

Un punto clave a considerar sobre el comportamiento de las personas es el valor y peso que tiene la moneda en el tiempo, de esta manera, es nesesario considerar la inflacion, y el aumento del precio en los distintos insumos del transporte.

\begin{table}[H]
    \centering
    \caption{Inflacion en Chile}
    \vspace{0.2cm}
    \begin{tabular}{|c|c|}
        \hline
        Año & Inflacion \\
        \hline
        2012 - 2017 & 18,2\% \\
        2012 - 2023 & 61,5\% \\
        \hline
    \end{tabular}
    \vspace{0.2cm}
    \\Fuente: \textbf{\cite{ipc}}
\end{table}

Donde, algunos de los insumos relacionados con el transporte que han aumentado de precio son la bencina y el pesaje de transporte publico (usaremos el metro) estos se encuentran en la tabla \textbf{\ref{Insumos}}. Las variaciones porcentuales fueron calculadas sobre el año 2012

Considerando los resultados obtenidos, se puede ajustar el parametro Xi para ajustar la regrecion a los distintos años:

\begin{table}[H]
    \centering
    \caption{Ajuste de parametros Xi}
    \vspace{0.2cm}
    \begin{tabular}{|c|c|c|}
        \hline
        Año & $X_i$ & Numero de Viajes \\
        \hline
        2012 & 5.730 & 1.84 \\
        2017 & 5.700 & 1.74 \\
        2023 & 5.725 & 1.58 \\
        \hline
    \end{tabular}
    \vspace{0.2cm}
    \\Fuente: Elaboracion propia
\end{table}


\subsection{Precisión de Resultados}
Analizando los datos presentados junto con las bases de datos, se estima que la precisión de los resultados obtenidos es fiable.
En primer lugar, las fuentes de donde se recopiló la información fueron extraídas directamente de las páginas web gubernamentales y municipales, 
así como también de los agentes encargados de análisis estadísticos sociales, como el INE.\\ 

En la sección 2.1 se analizaron datos de una muestra poblacional, mientras que en la sección 2.2 son de la población total en ese año.
Lo anterior se expresa en la diferencia de los parámetros $X_p$, $X_e$, $X_t$ y $X_j$. Para visualizar la precisión de mejor manera se generó la
siguiente tabla con valores estadísticos.\\

Para llamar la tabla de precision de resultados, se debe poner el siguiente comando \textbf{\ref{Precision_Resultados}}

De la tabla anterior se puede apreciar que los resultados obtenidos poseen una desviación estándar insignificante, por lo que se 
consideran representativos para cada coeficiente y su respectiva muestra.\\

Además, se realizó un análisis del estadístico t para 
comparar la media muestral con la media poblacional, teniendo como hipótesis cero que la media poblacional es menor a la media muestral.
Los resultados para cada rango etario demostraron que no se rechaza esa hipótesis, ya que \(\left| p_{value} \right| < \left| t_{value} \right|\), es decir, existe una correlación entre las medias, 
por lo que se puede afirmar que los resultados obtenidos son precisos y descriptivos.

\section{Anexos}

\subsection{Tablas Inflacion}

\begin{table}[H]
    \centering
    \caption{Aumento de indicadores economicos}
    \vspace{0.2cm}
    \begin{tabular}{|c|c|c|c|c|}
        \hline
        Año & UF (Pesos) & $\Delta$ \% UF & Valor Dolar & $\Delta$ \% Dolar\\
        \hline
        2012 & 22.296,19 & - & 501.34 & - \\
        2017 & 26.348,83 & 18.2\% & 661.19 & 31.88 \\
        2023 & 35.122,26 & 57.6\% & 826.34 & 64.82 \\
        \hline
    \end{tabular}
    \vspace{0.2cm}
    \\Fuente: \textbf{\cite{sii}}
\end{table}

\begin{table}[H]
    \centering
    \caption{Aumento de precio de insumos}
    \vspace{0.2cm}
    \begin{tabular}{|c|c|c|c|c|}
        \hline
        Año &  Bencina 95 (USD) & $\Delta$ \% Bencina & Metro (CLP) & $\Delta$ \% Metro\\
        \hline
        2012 &  1.56  & - & 610 & - \\
        2017 &  1.15 & -26.28 & 660 & 8.19\\
        2023 &  1.64 & 5.13 & 730 & 19.67\\
        \hline
    \end{tabular}
    \label{Insumos}
    \vspace{0.2cm}
    \\Fuente: \textbf{\cite{tradingeconomics}, \cite{bcentral}}
\end{table}

\subsection{Presicion de Datos}

\begin{table}[H]
    \caption{Análisis estadístico para la precisión de resultados.}
    \vspace{0.2cm}
    \centering
    \begin{tabular}{|c|c|c|c|}
    \hline
    \textbf{Coeficiente} & \textbf{Estadístico} & \textbf{Población} & \textbf{Muestra} \\ \hline
    $X_p$ & \multicolumn{1}{|c|}{\(\mu\)} & 0.178 & 0.220 \\ \cline{2-4}
    & \multicolumn{1}{|c|}{\(\sigma\)} & 0.010 & 0.053 \\ \cline{2-4}
    & \multicolumn{1}{|c|}{Valor p} & \multicolumn{2}{|c|}{1} \\ \cline{2-4}
    & \multicolumn{1}{|c|}{Valor t} & \multicolumn{2}{|c|}{-110.72} \\ \hline
    
    $X_e$ & \multicolumn{1}{|c|}{\(\mu\)} & 0.544 & 0.710 \\ \cline{2-4}
    & \multicolumn{1}{|c|}{\(\sigma\)} & 0.045 & 0.054 \\ \cline{2-4}
    & \multicolumn{1}{|c|}{Valor p} & \multicolumn{2}{|c|}{1} \\ \cline{2-4}
    & \multicolumn{1}{|c|}{Valor t} & \multicolumn{2}{|c|}{-99.52} \\ \hline
    
    $X_t$ & \multicolumn{1}{|c|}{\(\mu\)} & 1.539 & 1.910 \\ \cline{2-4}
    & \multicolumn{1}{|c|}{\(\sigma\)} & 0.074 & 0.117 \\ \cline{2-4}
    & \multicolumn{1}{|c|}{Valor p} & \multicolumn{2}{|c|}{1} \\ \cline{2-4}
    & \multicolumn{1}{|c|}{Valor t} & \multicolumn{2}{|c|}{-135.09} \\ \hline
    
    $X_j$ & \multicolumn{1}{|c|}{\(\mu\)} & 0.479 & 0.307 \\ \cline{2-4}
    & \multicolumn{1}{|c|}{\(\sigma\)} & 0.064 & 0.011 \\ \cline{2-4}
    & \multicolumn{1}{|c|}{Valor p} & \multicolumn{2}{|c|}{0} \\ \cline{2-4}
    & \multicolumn{1}{|c|}{Valor t} & \multicolumn{2}{|c|}{72.61} \\ \hline
    \end{tabular}
    \label{Precision_Resultados}
    \vspace{0.2cm}
    \\Fuente: Elaboración propia.
\end{table}

\subsection{Tablas Datos Encuesta EP}

\begin{table}[H]
    \centering
    \caption{Rango etario en Las Condes año 2012}
    \vspace{0.2cm}
    \begin{tabular}{|c|c|c|c|c|c|}
        \hline
        \textbf{Rango etario} & \textbf{Hombres} & \textbf{Mujeres} & \textbf{Total} & \textbf{\% Hombres} & \textbf{\% Mujeres} \\ \hline
        0-5 & 10484 & 9377 & 19861 & 52.79\% & 47.21\% \\ \hline
        6-22 & 33329 & 30366 & 63695 & 52.32\% & 47.67\% \\ \hline
        23-62 & 75938 & 82177 & 158115 & 48.02\% & 51.97\% \\ \hline
        63-80 & 17595 & 26095 & 43690 & 40.27\% & 59.72\% \\ \hline
    \end{tabular}
    \label{Cuadro 4}
    \vspace{0.2cm}
    \\Fuente: Elaboración propia.
\end{table}

\begin{table}[H]
    \centering
    \caption{Rango etario en Las Condes año 2017}
    \vspace{0.2cm}
    \begin{tabular}{|c|c|c|c|c|c|}
        \hline
        \textbf{Rango etario} & \textbf{Hombres} & \textbf{Mujeres} & \textbf{Total} & \textbf{\% Hombres} & \textbf{\% Mujeres} \\ \hline
        0-5 & 10672 & 9692 & 20364 & 52.41\% & 47.59\% \\ \hline
        6-22 & 31833 & 29028 & 60861 & 52.31\% & 47.69\% \\ \hline
        23-62 & 84879 & 88035 & 172914 & 49.08\% & 50.91\% \\ \hline
        63-80 & 21753 & 31816 & 53569 & 40.61\% & 59.39\% \\ \hline
    \end{tabular}
    \label{Cuadro 5}
    \vspace{0.2cm}
    \\Fuente: Elaboración propia.
\end{table}

\begin{table}[H]
    \centering
    \caption{Rango etario en Las Condes año 2023}
    \vspace{0.2cm}
    \begin{tabular}{|c|c|c|c|c|c|}
        \hline
        \textbf{Rango etario} & \textbf{Hombres} & \textbf{Mujeres} & \textbf{Total} & \textbf{\% Hombres} & \textbf{\% Mujeres} \\ \hline
        0-5 & 10769 & 9758 & 20527 & 52.46\% & 47.53\% \\ \hline
        6-22 & 31416 & 28678 & 60094 & 52.27\% & 47.72\% \\ \hline
        23-62 & 96289 & 97341 & 193630 & 49.72\% & 50.27\% \\ \hline
        63-80 & 27825 & 39107 & 66932 & 41.57\% & 58.42\% \\ \hline
    \end{tabular}
    \label{Cuadro 6}
    \vspace{0.2cm}
    \\Fuente: Elaboración propia.
\end{table}

\bibliographystyle{plainnat}  % Estilo de la bibliografía (puedes cambiarlo a otro si prefieres)
\bibliography{referencias.bib}  % Nombre del archivo .bib

\end{document}
